\section{Протокол и платформа}
\label{sec:protocol}

\subsection{Обзор альтернатив}
% ЧЕРНОВИК
SecureML~\cite{mohassel2017secureml}, SecureNN~\cite{wagh2019securenn},
Falcon~\cite{wagh2021falcon}, CrypTen~\cite{knott2021crypten}; обзоры
трёх- и четырёхсторонних систем~\cite{zapechnikov2022three,zapechnikov2022four}.
Большинство этих работ ориентировано на нейронные сети, где узким местом
являются умножения и аппроксимации~\cite{lapina2024}; для решающего дерева
узкое место~--- сравнения.

\subsection{Обоснование выбора протокола ABY3}
ABY3~\cite{mohassel2018aby3} использует реплицированное разделение секрета
на кольце $\mathbb{Z}_{2^{64}}$~\cite{araki2016} и переходы между
арифметическим, булевым и Yao-представлениями.

\subsection{Обоснование выбора платформы TF Encrypted}
TF Encrypted~\cite{dahl2018tfe} предоставляет Keras-подобный интерфейс и
реализацию ABY3.

\subsection{Параметры протокола}
Значение $x$ кодируется в фиксированной точке с $f$ дробными битами:
\begin{equation}
    \label{eq:fixed-point}
    \tilde{x} = \lfloor x \cdot 2^{f} \rceil \bmod 2^{64}.
\end{equation}

Для решающего дерева в фиксированной точке кодируются признаки объекта и
пороги в узлах; разрядность $f$ должна сохранять порядок значений, иначе
меняется результат сравнения.

\subsection{Вычисление решающего дерева в МКВ}
% ЧЕРНОВИК
Модель~--- решающее дерево классификации глубины $d$. В отличие от
нейронных сетей, аппроксимации нелинейных функций (сигмоиды, softmax) не
требуются: основная операция~--- конфиденциальное сравнение признака с
порогом в каждом узле. В ABY3 сравнение выполняется переходом от
арифметического представления к булевому и извлечением старшего
бита~\cite{mohassel2018aby3}.

Чтобы путь по дереву не раскрывал данные, вычисляются результаты
сравнений во всех $2^{d}-1$ внутренних узлах, после чего значение листа
выбирается конфиденциально (произведение битов сравнения вдоль пути).
Стоимость растёт экспоненциально по $d$, поэтому глубина дерева~---
ключевой параметр экспериментов.

В TF Encrypted нет готовой реализации решающего дерева: её предстоит
собрать из примитивов сравнения и умножения протокола ABY3. Это главный
технический риск, который проверяется при пробном запуске
(раздел~\ref{sec:trial}).
